%%%%%%%%%%%%%%%%%
% This is an sample CV template created using altacv.cls
% (v1.7.2, 28 August 2024) written by LianTze Lim (liantze@gmail.com). Compiles with pdfLaTeX, XeLaTeX and LuaLaTeX.
%
%% It may be distributed and/or modified under the
%% conditions of the LaTeX Project Public License, either version 1.3
%% of this license or (at your option) any later version.
%% The latest version of this license is in
%%    http://www.latex-project.org/lppl.txt
%% and version 1.3 or later is part of all distributions of LaTeX
%% version 2003/12/01 or later.
%%%%%%%%%%%%%%%%

%% Use the "normalphoto" option if you want a normal photo instead of cropped to a circle
% \documentclass[10pt,a4paper,normalphoto]{../_shared/altacv}

\documentclass[8pt,a4paper,ragged2e,withhyper]{../_shared/altacv}
%% AltaCV uses the fontawesome5 and simpleicons packages.
%% See http://texdoc.net/pkg/fontawesome5 and http://texdoc.net/pkg/simpleicons for full list of symbols.
%EXPERIENCE
%Senior researcher, CNES : LISA esa/NASA Space mission
% Change the page layout if you need to
\geometry{left=1.25cm,right=1.25cm,top=.5cm,bottom=1.cm,columnsep=1.2cm}

% The paracol package lets you typeset columns of text in parallel
\usepackage{paracol}
\usepackage{fontawesome5}
\usepackage{hyperref}

% Change the font if you want to, depending on whether
% you're using pdflatex or xelatex/lualatex
% WHEN COMPILING WITH XELATEX PLEASE USE
% xelatex -shell-escape -output-driver="xdvipdfmx -z 0" sample.tex
\iftutex 
  % If using xelatex or lualatex:
  \setmainfont{Roboto Slab}
  \setsansfont{Lato}
  \renewcommand{\familydefault}{\sfdefault}
\else
  % If using pdflatex:
  \usepackage[rm]{roboto}
  \usepackage[defaultsans]{lato}
  % \usepackage{sourcesanspro}
  \renewcommand{\familydefault}{\sfdefault}
\fi

% Change the colours if you want to
\definecolor{SlateGrey}{HTML}{2E2E2E}
\definecolor{LightGrey}{HTML}{666666}
\definecolor{DarkPastelRed}{HTML}{8AC0D9}%{450808}
\definecolor{PastelRed}{HTML}{00B0DA}%{8F0D0D}
\definecolor{GoldenEarth}{HTML}{B4AD0C}%{E7D192}
\colorlet{name}{black}
\colorlet{tagline}{PastelRed}
\colorlet{heading}{DarkPastelRed}
\colorlet{headingrule}{GoldenEarth}
\colorlet{subheading}{PastelRed}
\colorlet{accent}{PastelRed}
\colorlet{emphasis}{SlateGrey}
\colorlet{body}{LightGrey}

\definecolor{violet}{HTML}{5A2582}
\definecolor{rouge}{HTML}{F53B3B}
\definecolor{noir}{HTML}{00232C}
\definecolor{bleu}{HTML}{00B0DA}
\definecolor{rose}{HTML}{F831A0}
\definecolor{jaune}{HTML}{FFEC00}
\definecolor{vert}{HTML}{34AD6C}


% Change some fonts, if necessary
\renewcommand{\namefont}{\Huge\rmfamily\bfseries}
\renewcommand{\personalinfofont}{\footnotesize}
\renewcommand{\cvsectionfont}{\LARGE\rmfamily\bfseries}
\renewcommand{\cvsubsectionfont}{\large\bfseries}
%\newcommand{\orcid}[1]{\href{https://orcid.org/#1}{\includesvg[width=10pt]{orcid}}}

% Change the bullets for itemize and rating marker
% for \cvskill if you want to
\renewcommand{\cvItemMarker}{{\small\textbullet}}
\renewcommand{\cvRatingMarker}{\faCircle}
% ...and the markers for the date/location for \cvevent
% \renewcommand{\cvDateMarker}{\faCalendar*[regular]}
% \renewcommand{\cvLocationMarker}{\faMapMarker*}


% If your CV/résumé is in a language other than English,
% then you probably want to change these so that when you
% copy-paste from the PDF or run pdftotext, the location
% and date marker icons for \cvevent will paste as correct
% translations. For example Spanish:
% \renewcommand{\locationname}{Ubicación}
% \renewcommand{\datename}{Fecha}


%% Use (and optionally edit if necessary) this .tex if you
%% want to use an author-year reference style like APA(6)
%% for your publication list
% \input{../_shared/pubs-authoryear.tex}

%% Use (and optionally edit if necessary) this .tex if you
%% want an originally numerical reference style like IEEE
%% for your publication list
\input{../_shared/pubs-num.tex}

%% sample.bib contains your publications
\addbibresource{../_shared/sample.bib}
% \usepackage{academicons}\let\faOrcid\aiOrcid
\begin{document}
\name{BOILEAU Guillaume}
\tagline{Senior Python Engineer | Backend, APIs, Data \& AI Workflows}%Your Position or Tagline Here}
%% You can add multiple photos on the left or right
\photoL{2.8cm}{../_shared/moi}
%\photoR{2.5cm}{../_shared/Suitcase_High}%6892(1)}

\personalinfo{%
  % Not all of these are required!
  \email{guillaume.boileaupro@gmail.com}
  \phone{+33 6 59 94 85 84}
  \location{Alpes Maritimes, FRANCE}
  %\mailaddress{1 avenue des citronniers 06800 Cagnes sur Mer, France}
  %\homepage{https://www.researchgate.net/profile/Guillaume-Boileau-2}%\href{https://www.researchgate.net/profile/Guillaume-Boileau-2}{Guillaume-Boileau}}
 %\faLinkedin \href{https://www.linkedin.com/in/guillaume-boileau-phd-891b81265/}{in/guillaume-boileau}
     

  % \twitter{@twitterhandle}
  %\xtwitter{@x-handle}
  \linkedin{guillaume-boileau-phd-891b81265}
  %\github{your_id}
  \researchgate{Guillaume-Boileau-2}
  \orcid{0000-0002-3576-6968}
  
  %% You can add your own arbitrary detail with
  %% \printinfo{symbol}{detail}[optional hyperlink prefix]
  % \printinfo{\faPaw}{Hey ho!}[https://example.com/]

  %% Or you can declare your own field with
  %% \NewInfoFiled{fieldname}{symbol}[optional hyperlink prefix] and use it:
  % \NewInfoField{gitlab}{\faGitlab}[https://gitlab.com/]
  % \gitlab{your_id}
  %%
  %% For services and platforms like Mastodon where there isn't a
  %% straightforward relation between the user ID/nickname and the hyperlink,
  %% you can use \printinfo directly e.g.
  % \printinfo{\faMastodon}{@username@instace}[https://instance.url/@username]
  %% But if you absolutely want to create new dedicated info fields for
  %% such platforms, then use \NewInfoField* with a star:
  % \NewInfoField*{mastodon}{\faMastodon}
  %% then you can use \mastodon, with TWO arguments where the 2nd argument is
  %% the full hyperlink.
  % \mastodon{@username@instance}{https://instance.url/@username}
}
\vspace{-0.15cm}
\makecvheader
%% Depending on your tastes, you may want to make fonts of itemize environments slightly smaller

\vspace{-0.48cm}
\AtBeginEnvironment{itemize}{\small}

%% Set the left/right column width ratio to 6:4.
\columnratio{0.64}

% Start a 2-column paracol. Both the left and right columns will automatically
% break across pages if things get too long.
\begin{paracol}{2}
\cvsection{Experience}

\cvevent{\textbf{Software Engineer – Backend Services, System Interfaces \& Operational Reliability}}{VEV Platform Services France}{2025 -- 2026}{Valbonne, France}

\textbf{Context:} Development and integration of software services for EV charging infrastructure within a CPMS platform.

\textbf{Objective:} Deliver reliable backend services, operational data flows and robust system interfaces aligned with production constraints.

\begin{itemize}
\item \textbf{Backend service development:} Developed and maintained software services supporting operational monitoring and data exchange with charging stations.
\item \textbf{System interfaces:} Contributed to service integration involving software/hardware interactions and operational platform constraints.
\item \textbf{Service reliability:} Supported investigation and resolution of production issues affecting service continuity and data exchange consistency.
\item \textbf{Operational data flows:} Worked on structured data exchange services and service behaviour monitoring in production-related environments.
\item \textbf{Technical quality:} Contributed to maintainable implementations and structured follow-up of incidents and service improvements.
\end{itemize}

\divider

\cvevent{\textbf{Senior R\&D Engineer – Python Platforms, Data Workflows \& Technical Delivery}}{Laboratoire Lagrange, Observatoire de la Côte d'Azur, CNES}{2023 -- 2025}{Nice, France}

\textbf{Context:} Development and coordination of software and analysis activities for a large-scale international space mission involving complex data workflows.

\textbf{Objective:} Design and deliver robust Python-based tools and workflows for large-scale data processing, model integration and technical analysis.

\begin{itemize}
\item \textbf{Python platform development:} Developed a Python-based platform for large-scale model integration, comparison and analysis workflows. \textcolor{DarkPastelRed}{\href{https://gitlab.oca.eu/gboileau/lisa_l3_tools}{\bf GitLab:CATpyLISA}}
\item \textbf{Workflow engineering:} Structured modular analysis pipelines and processing workflows in a demanding collaborative environment.
\item \textbf{Data-intensive processing:} Worked on large-scale technical analysis and complex data processing scenarios requiring robustness and reproducibility.
\item \textbf{Software quality and maintainability:} Contributed to consistency checks, structured outputs, documentation and continuous improvement of workflows and methods.
\item \textbf{Technical investigation:} Analysed anomalies, investigated unexpected behaviours and supported reliable resolution of technical issues.
\item \textbf{Technical coordination:} Coordinated contributors and supported engineering activities across technical tasks.
\end{itemize}
\divider

\cvevent{\textbf{Data Scientist – Python Analysis Pipelines \& Performance Engineering}}{Universiteit Antwerpen, Belgium}{2021 -- 2023}{Antwerp, Belgium}

\textbf{Context:} Data analysis and performance studies for precision instruments in a high-requirement technical environment.

\textbf{Objective:} Develop robust Python workflows for performance analysis, technical investigation and improvement of system sensitivity.

\begin{itemize}
\item \textbf{Python workflows:} Developed structured pipelines for identification and characterization of performance-limiting effects.
\item \textbf{Data analysis:} Analysed complex datasets to investigate system behaviour and support improvement decisions.
\item \textbf{Technical robustness:} Worked in a context requiring reliability, reproducibility, traceability and structured reporting.
\item \textbf{Performance investigation:} Studied environmental and technical factors affecting measurement quality and system-level performance.
\end{itemize}
\divider

\cvevent{\textbf{Junior R\&D Engineer – Python Modeling, Data Analysis \& Technical Studies}}{ARTEMIS, Observatoire de la Côte d'Azur, Nice, France}{2018--2021}{Nice, France}

\textbf{Context:} Engineering activities involving complex data analysis, simulation and technical studies in a collaborative environment.

\textbf{Objective:} Develop robust Python-based methods and workflows for large-scale data analysis and technical investigations.

\begin{itemize}
\item \textbf{Data analysis:} Developed methods for the analysis of complex low-SNR datasets.
\item \textbf{Modeling workflows:} Built simulation approaches for large-scale datasets and complex technical environments.
\item \textbf{Technical studies:} Contributed to collaborative investigations and structured analysis activities.
\item \textbf{Method development:} Participated in the development of analysis tools supporting shared technical objectives.
\end{itemize}




\switchcolumn

\iffalse
\cvsection{My Life Philosophy}

\begin{quote}
``Life is like riding a bicycle. To keep your balance, you must keep moving.''
\end{quote}


\cvsection{Summary and Strengths}

Experienced astrophysicist with advanced skills in statistical analysis, Bayesian modeling, and machine learning. Proven ability to lead teams in international collaborations. Strong scientific background demonstrated by impactful publications and presentations at major conferences. Skilled in managing large-scale scientific projects and solving complex problems.
\fi

\cvsection{Summary and Strengths}

Senior engineer with strong experience in \textbf{Python development}, \textbf{data-intensive workflows}, \textbf{backend-oriented engineering} and \textbf{complex technical environments}. Used to designing reliable software components, structuring analysis and processing pipelines, and delivering robust technical solutions in demanding contexts.

Strong background in \textbf{data and AI-related workflows}, \textbf{technical investigation}, \textbf{system integration} and \textbf{software quality}. Comfortable working on modular codebases, service interactions, structured data processing and continuous improvement of technical workflows.

Hands-on experience with \textbf{Python}, \textbf{SQL}, \textbf{REST-oriented services}, \textbf{Docker}, and modern software development practices in environments requiring rigor, maintainability, traceability and collaboration.

\cvsection{Team Management}

\begin{itemize}
  \item Supervision of \textbf{3 Master’s thesis interns (M2)} during engineering and development phases.
  \item Coordination of technical activities involving \textbf{research engineers and technicians} on a complex data processing pipeline.
  \item Experience in organizing tasks, aligning contributors, supporting delivery and maintaining structured communication across collaborative environments.
\end{itemize}

\cvsection{Languages}
%\cvskill{French}{5}
%\divider
\cvskill{English}{4}
\divider

\cvskill{Spanish}{2}
\divider

%\cvskill{German}{3.5} %% Supports X.5 values.

%% Yeah I didn't spend too much time making all the
%% spacing consistent... sorry. Use \smallskip, \medskip,
%% \bigskip, \vspace etc to make adjustments.
\medskip

\cvsection{Education}
\cvevent{PhD in Physics}{Université Côte d'Azur}{2018 -- 2021}{Nice, France}
\divider


\cvevent{Selective Graduate Program in Fundamental Physics (Magistère)}{Université Paris-Saclay}{2016 -- 2018}{Orsay, France}
%\textit{with distinction}

\divider
\cvevent{MSc in Astronomy and Astrophysics}{Observatoire de Paris / Université Paris-Saclay}{2017 -- 2018}{Paris, France}
%\textit{with distinction}
\divider
\cvevent{BSc in Fundamental Physics}{Université Paris-Saclay}{2015 -- 2016}{Orsay, France}
\cvsection{Professional Training}

\cvevent{Supervised Machine Learning (Regression \& Classification)}{DeepLearning.AI – Andrew Ng}{2020}{Coursera}
%\textit{with distinction}

%\divider

%\cvevent{CPGE Classe préparatoire aux grandes écoles MPSI/MP (Maths-Physics)}{Lycée  Dessaigne}{2013-2015}{Blois, France}
%\textit{with distinction and merit}

%\divider

%\cvevent{High School Diploma, Baccalauréat  S }{Lycée  Augustin  Thierry}{2013}{Blois, France}

%\cventry{2013-2015}{CPGE Classe préparatoire aux grandes écoles MPSI/MP (Maths-Physics)}{Lycée  Dessaigne, Blois, France}{%(a  two  years preparatory  course  that  is  a  prerequisite  for  entry  into  one  of  the Grandes Ecoles)
%}{}{\textit{with distinction and merit}}

%\cventry{2013}{High School Diploma, Baccalauréat  S %(high  school diploma focusing on Mathematics and Physics) 
%}{Lycée  Augustin  Thierry}{Blois, France}{}{\textit{with distinction }}


\end{paracol}

\newpage
\cvsection{Projects}

\cvevent{LISA Consortium}{Large-scale scientific software and data environment}{2018--now}{}

Contribution to complex software and data workflows in an international environment involving large-scale analysis, technical coordination and robust Python-based tooling.

\textbf{Role:} Development of Python workflows, analysis tools and structured technical methods for complex processing environments.

\divider

\cvevent{Einstein Telescope and LIGO-Virgo-KAGRA Collaborations}{Collaborative data and software environments}{2021--now}{}

Contribution to collaborative software and data analysis developments in demanding technical environments.

\textbf{Role:} Development of analysis tools, contribution to technical workflows and support to structured software methods.
%\medskip

\cvsubsection{Core skills}
{\small
\cvtag{Python Development}
\cvtag{Backend Engineering}
\cvtag{API Design}
\cvtag{REST Services}
\cvtag{Data Workflows}
\cvtag{AI Workflows}
\cvtag{Software Architecture}
\cvtag{Modular Design}
\cvtag{Technical Investigation}
\cvtag{System Integration}
\cvtag{Code Quality}
\cvtag{Maintainability}
\cvtag{Performance Analysis}
\cvtag{Scalable Processing}
\cvtag{Continuous Improvement}
\cvtag{Technical Documentation}
}
\\
\divider\smallskip
\vspace{-.3cm}

\cvsubsection{Technical skills}
{\small
\cvtag{Python}
\cvtag{SQL}
\cvtag{Relational Databases}
\cvtag{MongoDB}
\cvtag{CI/CD}
\cvtag{REST API}
\cvtag{Docker}
\cvtag{Linux}
\cvtag{Git}
\cvtag{TypeScript}
\cvtag{Node.js}
\cvtag{Pandas}
\cvtag{PyTorch}
\cvtag{Data Analysis}
\cvtag{Technical Reporting}
\cvtag{Workflow Automation}
\cvtag{Software Engineering}
}
\\
\divider\smallskip
\vspace{-.3cm}

\cvsubsection{Soft skills}
{\small
\cvtag{Autonomy}
\cvtag{Analytical Thinking}
\cvtag{Problem Solving}
\cvtag{Collaboration}
\cvtag{Technical Communication}
\cvtag{Execution Rigor}
\cvtag{Adaptability}
\cvtag{Reliability}
\cvtag{Attention to Detail}
\cvtag{Continuous Learning}
\cvtag{Proactiveness}
\cvtag{Structured Communication}
\par}

\cvtag{\textit{Applications: backend services, data platforms, AI workflows, complex software systems}}
\end{document}
